\documentclass[a4paper,12pt]{article}

%%% Alapdokumentum beállítások %%%
\usepackage[hungarian]{babel}     % Magyar nyelvi támogatás
\usepackage[T1]{fontenc}          % T1 betűtípus kódolás (az utf8 inputenc ma már alapértelmezett)
\usepackage{geometry}             % Oldalformázás
\geometry{a4paper, margin=1.5in}  % Papírméret és margók

%%% Szövegformázás és oldalelrendezés (JAVÍTOTT) %%%
\usepackage{microtype}            % Intelligens sorkizárás, szebb szövegkép
\usepackage{parskip}              % Bekezdések közötti térköz (behúzás helyett)
\usepackage{enumitem}             % Lista testreszabáshoz
\usepackage{multicol}             % Többszörös oszlopokhoz
\setlength{\emergencystretch}{3em}% Segít a nehezen törhető soroknál (a \sloppy helyett)
% Az elválasztás tiltása (\hyphenpenalty) kivezetve, mert magyar nyelvben tönkreteszi a sorkizárást!

%%% Matematikai és tudományos csomagok %%%
\usepackage{amsmath, amssymb, amsfonts, amsthm} % Matematikai csomagok egyben
\usepackage{stmaryrd}                           % Speciális matematikai szimbólumok
\usepackage[version=4]{mhchem}                  % Kémiai képletekhez
\usepackage{esvect}                             % Vektorokhoz (vastag nyilak)
\usepackage{cancel}                             % Matematikai áthúzáshoz

%%% Színek és Grafikai elemek (JAVÍTOTT) %%%
\usepackage{xcolor}                             % Színek kezelése (FONTOS!)
\definecolor{bvred}{rgb}{0.8, 0.0, 0.0}         % Piros szín definiálása

\usepackage{graphicx}                           % Képek beszúrásához
\usepackage{wrapfig}                            % Képek lebegtetéséhez
\usepackage{tikz}                               % Vektoros ábrákhoz
\usetikzlibrary{calc, intersections, decorations.pathreplacing}
\usepackage{pgfplots}
\pgfplotsset{compat=1.18}                       % PGFPlots kompatibilitás (FONTOS!)
\usepackage{caption}                            % Ábrafeliratok testreszabásához
\usepackage{float}                              % Ábrák pozicionálásához

%%% Táblázatok és tömbök %%%
\usepackage{array}                       % Speciális oszloptípusok
\setlength{\arrayrulewidth}{0.8pt}       % Táblázatvonalak vastagsága
\setlength{\tabcolsep}{6pt}              % Oszlopok közötti térköz
\renewcommand{\arraystretch}{1.8}        % Sorok magassága

%%% Egyedi parancsok és környezetek (JAVÍTOTT) %%%
\newcommand{\textbfr}[1]{\textbf{\textcolor{bvred}{#1}}} % Piros félkövér szöveg

\usepackage{etoolbox} % LaTeX programozási eszközök
\newrobustcmd*{\mytriangle}[1]{%                  % Háromszög jelölés
    \tikz{\filldraw[draw=#1,fill=#1,rotate=-90] (0,0)--(0.2cm,0)--(0.1cm,0.2cm)--cycle;}}

% amsthm környezetek letisztult beállítása
\theoremstyle{definition} % Ezzel a normál (nem dőlt) betűtípust használja a környezeteken belül
\newtheorem*{definition}{Definíció}
\newtheorem*{theorem}{Tétel}
\newtheorem*{problem}{Feladat}
\newtheorem*{example}{Példa}
\newtheorem*{remark}{\mytriangle{bvred} Megjegyzés}

%%% Hiperhivatkozások és PDF-beállítások %%%

%%% Megjelenítéshez automatikusan kiegészített csomagok %%%
\usepackage[utf8]{inputenc}
\usepackage{tcolorbox}
\tcbuselibrary{theorems,skins,breakable}
\usepackage{booktabs}

%%% Megjelenítéshez kiegészített dobozok és makrók %%%
\newtcolorbox{definiciodoboz}[1][]{colback=blue!5!white,colframe=blue!75!black,fonttitle=\bfseries,title=Definíció: #1}
\newtcolorbox{teteldoboz}[1][]{colback=green!10!white,colframe=green!75!black,fonttitle=\bfseries,title=Tétel: #1}
\newtcolorbox{torvenydoboz}[1][]{colback=purple!5!white,colframe=purple!75!black,fonttitle=\bfseries,title=Fizikai törvény: #1}

\usepackage{hyperref}                    % Mindig utolsóként (vagy majdnem utolsóként) érdemes betölteni

%%% Geometriai jelölések %%%
\def\dotMarkRightAngle[size=#1](#2,#3,#4){%
 \draw ($(#3)!#1!(#2)$) -- 
       ($($(#3)!#1!(#2)$)!#1!90:(#2)$) --
       ($(#3)!#1!(#4)$);
 \path (#3)--node[circle,fill,inner sep=.5pt]{} ($($(#3)!#1!(#2)$)!#1!90:(#2)$);
}

% AI által javasolt kiegészítések:
\newtcolorbox{torvenydoboz}[1][]{colback=purple!5!white,colframe5!purple!75!black,fonttitle=\bfseries,title=Fizikai törvény: #1}

\title{Ismétlés nélküli kombináció}
\author{Készítette: Varga Dávid}

\begin{document}

\maketitle

\section{Mikor használjuk?}
Képzeld el, hogy a gyümölcsöskosárba kell válogatnod gyümölcsöket, vagy éppen a lottón kell megjelölnöd néhány számot. Az **ismétlés nélküli kombinációt** akkor hívjuk segítségül, amikor egy nagyobb halmazból kell kiválasztanunk egy adott darabszámú kisebb csoportot úgy, hogy **a sorrend egyáltalán nem számít**, és elemeket **nem halmozhatunk** (azaz mindent legfeljebb egyszer választhatunk ki). 

A lényeg tehát két pontban sűríthető össze: 
1. Mindegy, hogy a kiválasztott elemeket milyen sorrendben kapjuk meg (ha az alma, körte, banán hármast húzzuk ki, az teljesen ugyanaz, mintha körte, banán, alma lenne a sorrend).
2. Nincs visszatevés, mindegyik elem csak egyszer szerepelhet a pakliban.

\section{Hogyan működik?}
Nézzünk meg egy szemléletes példát a logika megértéséhez! Tegyük fel, hogy van $5$ féle kiváló csokoládénk a boltban ($A, B, C, D, E$), és nekünk pontosan **$3$ darazat** kell választanunk a csomagunkba. 

Ha a **sorrend számítana** (vagyis az elsőként választott csoki különleges státuszt kapna, mondjuk ő lenne a díszdoboz közepén), akkor így gondolkodhatnánk:
Az első helyre választhatunk $5$félét, a másodikra $4$-et, a harmadikra pedig $3$-at. Ez összesen $5 \cdot 4 \cdot 3 = 60$ lehetőséget jelentene (ez a \textbf{variáció}).

De figyeljünk oda! A mi esetünkben a csokik egy sima zacskóba kerülnek, ott pedig nincsenek "helyek". Ha a zacskóba bekerül az $A$, a $B$ és a $C$ csoki, az pont ugyanaz a végeredmény, mintha a $C$, az $A$ és a $B$ csokoládét pakoltuk volna be. Hányféleképpen tudjuk sorba rendezni ezt a $3$ kiválasztott csokit? Pontosan $3!$ ( azaz $3 \cdot 2 \cdot 1 = 6$) féleképpen. 

Mivel ez a $6$-féle sorrend a mi szempontunkból mind **ugyanazt a kombinációt** jelenti, a korábban kapott $60$-as számot túl sokat számoltunk! Ahhoz, hogy ezt a felesleges duplikációt kiszűrjük, a lehetőségek számát egyszerűen el kell osztanunk a sorrendek számával, vagyis $3!$-mal:
$$\frac{5 \cdot 4 \cdot 3}{3!} = \frac{60}{6} = 10.$$
Ez a bizonyos "osztás" a titka annak, hogy a sorrendet érvénytelenítsük.

\section{Általános képlet és speciális esetek}
Ha az előbbi logikát átültetjük a matematika nyelvérre: van összesen $n$ darab különböző elemünk, és ezek közül szeretnénk kiválasztani $k$ darazat (ahol $k \le n$). Az ismétlés nélküli kombinációk száma a következő képlettel adható meg:

\begin{teteldoboz}{Ismétlés nélküli kombináció képlete}
$$\binom{n}{k} = \frac{n!}{k!(n-k)!}$$
Ezt úgy olvassuk ki: „$n$ alatt a $k$”.
\end{teteldoboz}

\begin{remark}
Gyakori buktatók és speciális esetek, amikre érdemes figyelni:
\begin{itemize}
    \item \textbf{A sorrend csapdája:} Ha a feladatban a sorrendnek, rangnak, beosztásnak szerepe van (pl. elnök, alelnök, titkár választás), akkor \textit{nem} kombinációt, hanem variációt használunk!
    \item \textbf{Határesetek:} $\binom{n}{0} = 1$ (sehogy sem választani 1féleképpen lehet) és $\binom{n}{n} = 1$ (az összes elemet kiválasztani szintén 1féleképpen lehet).
    \item \textbf{Szimmetria-tulajdonság:} $\binom{n}{k} = \binom{n}{n-k}$. Ha pl. $10$ emberből $8$-at akarunk kiválasztani, az ugyanannyi lehetőség, mintha azt néznénk, hogy melyik $2$ embert \textit{hagyjuk otthon}.
\end{itemize}
\end{remark}

\section{Dolgozat- / Vizsga-gyorsteszt (Döntési fa)}
Amikor olvasod a feladat szövegét, tedd fel magadnak az alábbi három kérdést sorban. Ha mindegyikre az alábbi válasz adódik, akkor azonnal tudod, hogy ismétlés nélküli kombinációt kell írnod:

\begin{enumerate}[label=\textbfr{\arabic*}.]
    \item \textbf{Van-e visszatevés / ismétlődés?} $\rightarrow$ Nincs (minden elem egyedi, mindent csak egyszer választhatunk).
    \item \textbf{Számít-e a kiválasztott elemek sorrendje?} $\rightarrow$ Nem számít (csak a csoport összetétele a lényeg, nincs első, második stb. helyezett).
    \item \textbf{Az összes rendelkezésre álló elemet felhasználjuk?} $\rightarrow$ Nem, csak egy részüket ($k$ darazat $n$-ből).
\end{enumerate}

\section{Tipikus és trükkös feladatok levezetése}

\begin{problem}[Kártyalapok húzása]
Egy standard magyar kártyapakliból ($32$ lapos) egyszerre kihúzunk $4$ lapot. Hányféle különböző 4 lapos kéz jöhet létre?
\end{problem}

\textbf{Gondolatmenet és lépések:}
1. \textit{Mit látunk?} Van $n = 32$ darab különböző elemünk (a kártyák), és ebből választunk $k = 4$ darazat.
2. \textit{Miért így gondolkodunk?} Mivel a kártyákat egyszerre húzzuk ki (vagy ha sorban is húzzuk, de a végén csak a kezünkben lévő lapok halmaza számít), a sorrend teljesen mindegy. Egy kártyapakliból egy lap csak egyszer szerepel, így ismétlés sincs. Ez tiszta kombináció.
3. \textit{A számítás:}
$$\binom{32}{4} = \frac{32!}{4!(32-4)!} = \frac{32!}{4! \cdot 28!} = \frac{32 \cdot 31 \cdot 30 \cdot 29}{4 \cdot 3 \cdot 2 \cdot 1}$$
4. \textit{Hogyan rontják el gyakran?} Sok diák megijed a nagy faktoriálisoktól, vagy megpróbálja kiszorozni a $32!$-t a zsebszámológépén, ami túlcsordulást (overflow hibát) okoz. Mindig egyszerűsítsünk először a nevezővel!
5. \textit{Hibátlan megoldás:} 
A nevezőben lévő $4 \cdot 3 \cdot 2 \cdot 1 = 24$. Ezzel leosztva a számlálót: a $30$-at elosztjuk $3 \cdot 10$-zel vagy hasonlóan, kapjuk:
$$\frac{32 \cdot 31 \cdot 30 \cdot 29}{24} = 4 \cdot 31 \cdot 5 \cdot 29 = 35\,960.$$
Összesen $35\,960$-féleképpen húzhatunk ki $4$ lapot.

\vspace{1em}

\begin{problem}[Csapatépítés feltételekkel (Trükkös)]
Egy 10 fős baráti társaságból (4 fiú és 6 lány) szeretnénk egy 4 fős bizottságot alakítani. Hányféleképpen tehetjük ezt meg, ha a bizottságban \textbf{pontosan 2 fiúnak} kell lennie?
\end{problem}

\textbf{Gondolatmenet és lépések:}
1. \textit{Mit látunk?} Komplexebb feladat, feltétellel. A társaság két különálló részre oszlik: fiúkra és lányokra.
2. \textit{Miért így gondolkodunk?} Ha a bizottság 4 fős, és pontosan 2 fiút kell tartalmaznia, akkor logikus, hogy a maradék $4 - 2 = 2$ helyre lányok fognak kerülni. A fiúk kiválasztása független a lányok kiválasztásától, de a végén a lehetőségeket össze kell szoroznunk (\textbf{Kombinatorikus alapelv}: és-kapcsolat esetén szorzás).
3. \textit{A részfeladatok:}
   - Hányféleképpen tudunk kiválasztani 2 fiút a 4 fiúból? Ez egy ismétlés nélküli kombináció: $\binom{4}{2}$.
   - Hányféleképpen tudunk kiválasztani 2 lányt a 6 lányból? Ez szintén kombináció: $\binom{6}{2}$.
4. \textit{A számítás:}
   $$\binom{4}{2} = \frac{4 \cdot 3}{2 \cdot 1} = 6$$
   $$\binom{6}{2} = \frac{6 \cdot 5}{2 \cdot 1} = 15$$
5. \textit{Hogyan rontják el gyakran?} A diákok gyakran összeadják a két lehetőséget ($6 + 15 = 21$) a szorzás helyett. Emlékezzünk: a bizottsághoz \textit{szükségünk van} a fiúkra \textbf{és} a lányokra is egyszerre, ezért a kombinációk számát össze kell szorozni!
6. \textit{Hibátlan megoldás:}
   $$6 \cdot 15 = 90.$$
   Tehát $90$-féleképpen állíthatjuk össze a kívánt bizottságot.

\end{document}